\section{Polynomial estimates for the odd-prime positive eigenspace}
\label{sec:rh-plus-nonzero}

We derive the individual estimate \texttt{Height\allowbreak{}And\allowbreak{}Vanishing.\allowbreak{}Count\allowbreak{}Vs}
on the positive-eigenspace side of the strict window. In
\texttt{Joint\allowbreak{}Win\allowbreak{}Of\allowbreak{}Counts.\allowbreak{}point\allowbreak{}HV\_\allowbreak{}of\_\allowbreak{}counts}, this estimate is combined
with counts II, IV and the other branch of V to give \texttt{Sum\allowbreak{}In}.
The result concerns the twist \(e=1\), where the source gives an upper
bound for the length. We also derive the estimate used outside the
exceptional set for arbitrary positive twists in case I.

\subsection{The conductor condition bounds the degree}

Fix a primitive target \(\chi\) of conductor \(m\) and order \(o\),
and write \(M=\max B_\chi\), which exists since \(3\in B_\chi\).
Take any host of character modulus \(N\) satisfying the size conditions
of Section~\ref{sec:arithmetic-translation}, an odd prime
\(\ell\in B_\chi\), any admissible coefficient ring \(\mathcal O\), and
\(\psi\in\mathcal X_{\ell,1}\). The cutoff is \(t>1\), so \(e=1\) is
in the range. Write \(\psi_2\) for the \(2\)-component of \(\psi\)
at the level and \(c=v_2(\operatorname{cond}\psi_2)\).
The conductor of \(\psi_2\) is \(2^c\).

Remove the \(\ell\)-power-order part of \(\psi\), then its component
at two, to obtain \(\theta\); let \(a\) denote its order. Put
\[
  f=\operatorname{ord}_{\operatorname{lcm}(4,a)}(\ell),\qquad
  w=v_2(\ell^f-1),\qquad R=16o(\ell m)^2\,4^w.
\]
The hypothesis for this branch is \(2^c<R\). The host character
comparison \texttt{host\allowbreak{}Th1} gives \(a\mid2o\) and \(\ell\nmid a\).
The modulus \(d=\operatorname{lcm}(4,a)\) divides \(4o\) and is
prime to \(\ell\). Thus \(f\ge1\), \(f\mid\varphi(d)\), and
\(\varphi(d)\mid\varphi(4o)\). The last divisibility follows directly
on comparing the factors \(p^{v-1}(p-1)\) in Euler's product.
It follows that \(2^w\le\ell^f-1<\ell^{\varphi(4o)}\), giving
\[
        R\le 16o(\ell m)^2\ell^{2\varphi(4o)}.
\]
The source's component calculation \texttt{host\allowbreak{}Comp} bounds the order
of every component of \(\psi\) at an odd prime dividing \(N\) by
divisibility into \(2o\). The order of \(\psi_2\) divides
\(\varphi(2^c)\), since it factors through the units of its conductor,
and \(\varphi(2^c)\mid 2^c\), including \(c=0\).
The Chinese remainder decomposition now gives \(\psi^{2o\,2^c}=1\).
Thus \texttt{Pow\allowbreak{}Totient\allowbreak{}Conductor.\allowbreak{}vs\allowbreak{}Ord} gives
\[
       \operatorname{ord}(\psi)\,\ell\le
       D_0,\qquad
       D_0=2o\bigl(16o(Mm)^2 M^{2\varphi(4o)}\bigr)M .
\]
This number depends only on the target, even when \(N\) varies.

Inside \(\mathbb Q(\mu_{\ell N})\) take the fixed field \(L\) of
the kernel of the map from the units modulo \(\ell N\) given by
\[
            u\longmapsto \bigl(\psi(u\bmod N),u\bmod\ell\bigr).
\]
Then \(\mu_\ell\subset L\), \(L/\mathbb Q\) is abelian, and \(\psi\)
is trivial on its absolute Galois group. Its finite image in the first
factor is cyclic of order \(\operatorname{ord}(\psi)\), since it is a
finite subgroup of a field's multiplicative group. The second factor
has order \(\ell-1\). Consequently \(n=[L:\mathbb Q]\le
\operatorname{ord}(\psi)(\ell-1)\le D_0\). This is the construction in
\texttt{Vs\allowbreak{}Field.\allowbreak{}vs\allowbreak{}Field}. Put \(x=\ell N\) and
\(\Delta=|\operatorname{disc}L|\). The discriminant tower identity
gives \(\Delta^{1/n}\le
|\operatorname{disc}\mathbb Q(\mu_x)|^{1/\varphi(x)}\).
The cyclotomic discriminant formula bounds the right side by \(x\);
explicitly it is \(x\prod_{p\mid x}p^{-1/(p-1)}\).
Thus \(\log\Delta\le n\log x\), as in \texttt{Win.\allowbreak{}disc}.

\subsection{Descended length and the numerical comparison}

The arithmetic input \texttt{Lemma\allowbreak{}K} is a size assertion for \(H^2\).
At twist \(e\) over a totally complex field \(L\) containing the
specified roots of unity, and for every choice of its local data,
it gives a finite \(H^2\) for the compact integral cyclotomic twist
\(\mathbb Z_\ell(e+1)\) in the source's Galois problem, and
\[
 \log\# H^2 \le
 \exp\big((2en)^{C_K}\log(2\Delta)\big).
\]
In the present case \(L\) is totally complex since it contains a
primitive odd-prime root of unity. Write \(K=\max(0,\lceil C_K\rceil)\)
and \(A=(2D_0)^K(D_0+1)\). Replacing \(C_K\) by \(K\) increases
the right side. At \(e=1\), write \(k\) for the \(\mathbb Z_\ell\)-length
of this finite group, so its cardinality is \(\ell^k\). Because
\(\log\ell\ge\log2>1/2\),
\[
      k\le2\exp\big((2n)^K\log(2\Delta)\big)\le2x^A.
\]
For the second inequality use \(x\ge2\) and \(1\le n\le D_0\) to get
\(\log(2\Delta)\le\log x+n\log x\le(D_0+1)\log x\).
The exponent of \(x\) is therefore independent of the host and
coefficient choices.

We record the arithmetic comparisons used to descend this estimate.
The global-to-local inequality \texttt{Prop\allowbreak{}A.\allowbreak{}sha\allowbreak{}Le\allowbreak{}V} uses
Poitou--Tate duality, local finiteness and local Tate duality to
bound the strict \(\mathbb Z_\ell\)-length over \(L\) by \(k\).
The character restriction bound \texttt{Chr.\allowbreak{}up\allowbreak{}V} allows one error
from local degree-zero invariants for each generator of
\(\operatorname{Gal}(L/\mathbb Q)\). Take the full group as generating
set; it has \(n\) elements. The Galois character of \(\psi\) factors
through the units modulo \(N\), so
\texttt{H0\_\allowbreak{}local\_\allowbreak{}card\_\allowbreak{}bounded\_\allowbreak{}of\_\allowbreak{}cls\_\allowbreak{}ker\_\allowbreak{}le\_\allowbreak{}ker} applies.
With \(r_\ell=\operatorname{length}_{\mathcal O}(\mathcal O/\ell)\)
the resulting inequalities are
\[
    b_{\ell,1}(\psi)\ \le\ (k+n\nu)\,r_\ell,\qquad
    \nu=
    \begin{cases}
       \max(v_\ell(N),1),&\ell\mid N\ \text{or}\ \ell-1\mid 1,\\
       0,&\text{otherwise}.
    \end{cases}
\]
The displayed local constant specializes the definition
\texttt{LData\allowbreak{}Card\allowbreak{}Le\allowbreak{}Finrank.\allowbreak{}nu\allowbreak{}Bound} to odd \(\ell\) and \(e=1\).
It is at most \(2N\), since \(v_\ell(N)\le N\) and \(N\ge2\).
The displayed inequality also uses monotonicity when the strict
decomposition conditions at the primes dividing \(N\) are imposed.
All three comparisons here are separate arithmetic inputs.

The ramification index \(r_\ell\) is a positive finite integer,
so \(b_{\ell,1}(\psi)\) is finite and division gives
\(v_{\ell,1}(\psi)\le2x^A+2D_0 N\).
Since \(0\le\log\ell\le\ell\le x\) and \(N\ell=x\),
\[
\begin{aligned}
  (\log\ell)v_{\ell,1}(\psi)
   &\le 2x^A\ell+2D_0 x\\
   &\le 2(D_0+1)x^{A+1}\\
   &\le x^{A+D_0+3}.
\end{aligned}
\]
Indeed \(D_0+1\le2^{D_0}\) and \(x\ge2\), so
\(2(D_0+1)\le x^{D_0+2}\). Hence the asserted count holds with the
single exponent \((2D_0)^K(D_0+1)+D_0+3\). The two alternatives for
the conductor threshold are what permit the source to use this
degree bound for every occurrence of the strict inequality \(2^c<R\).

\subsection{The uniform case I estimate over a finite set of primes}

We now follow \texttt{Unitary\allowbreak{}Three\allowbreak{}Groups.\allowbreak{}Count\allowbreak{}I}. In the actual
window argument, \texttt{Linked\allowbreak{}Unl\allowbreak{}E.\allowbreak{}count\_\allowbreak{}i\_\allowbreak{}of\_\allowbreak{}printed\_\allowbreak{}bound\_\allowbreak{}k\_\allowbreak{}kato\_\allowbreak{}automorphy}
obtains its orbit estimate by applying \texttt{U3.\allowbreak{}cor\allowbreak{}GK} to
\texttt{prop\allowbreak{}GK\_\allowbreak{}of\_\allowbreak{}number\allowbreak{}U3\allowbreak{}U\_\allowbreak{}of\_\allowbreak{}one\allowbreak{}Char\allowbreak{}U}. Apply the orbit estimate
of Section~\ref{sec:arithmetic-translation} with
\(F=\mathbb Q(\sqrt{-p_K})\), and let \(c_5,A\) be its constants for
this field. Choose primes \(r_1>30mp_K\) and \(r_2>r_1\), and choose
\(k\ge0\) with \(2c_5(\log r_1+\log r_2)<16^k\). The source
claims the exponent \(k+1+A\) in count I for this target.

Indeed take a host for which each \(r_i\) is smaller than both
host primes \(q_1,d\), a twist \(1\le e<t\), an arbitrary family of
coefficient rings, and a finite set \(S\) of odd primes in case I.
The host-level facts give \(16\mid N\) and every prime divisor of
\(N\) equals two, divides \(m\), equals \(q_1\), or equals \(d\).
Furthermore \(|\operatorname{disc}F|=p_K\). Thus neither \(r_i\)
divides \(30N|\operatorname{disc}F|\). The test primes are eligible
simultaneously for every host satisfying these inequalities.

Let \(G_N=(\mathbb Z/N\mathbb Z)^\times\) and form the finite set
\(\mathcal H\) of kernels of characters in
\(\bigcup_{\ell\in S}\mathcal X_{\ell,e}\), taking kernels in \(G_N\).
For \(H\in\mathcal H\), let
\[
  S_H=\{\ell\in S:\ \ker\psi=H
           \text{ for some }\psi\in\mathcal X_{\ell,e}\}.
\]
For every \(\ell\in S_H\), the host lemma \texttt{gd\allowbreak{}Of\allowbreak{}Host} makes that
character a witness for the good-data predicate in the orbit estimate.
In particular every length occurring in the desired sum is finite.
Partition the character sum at each prime by its kernel, and enlarge
each nonempty part to the entire orbit of characters with that kernel.
The summands are nonnegative, and \(\log\ell\) is positive, so this
operation gives an upper bound. Applying the orbit estimate to each
\(H\) with its own \(S_H\) now gives
\[
 \sum_{\ell\in S}(\log\ell)\sum_{\psi\in\mathcal X_{\ell,e}}
       v_{\ell,e}(\psi)
 \le
 2c_5(\log r_1+\log r_2)\{N(e+1)\}^A
        \sum_{H\in\mathcal H}\varphi([G_N:H]).
\]
For each such \(H\), the quotient \(G_N/H\) is finite cyclic, being
isomorphic to the image of a multiplicative character in a field.
If its order is \(h\), the faithful complex characters on this
quotient number \(\varphi(h)\), by sending a generator to a primitive
\(h\)-th root of unity. Inflating them to \(G_N\) gives characters
with kernel exactly \(H\). Distinct \(H\)'s give disjoint sets.
The finite abelian group \(G_N\) has \(|G_N|\) complex characters:
decompose it as a direct product of cyclic groups and count the
choices on each factor. Therefore
\(\sum_{H\in\mathcal H}\varphi([G_N:H])\le |G_N|=\varphi(N)\le N\).
This proves the counting step \texttt{orb\allowbreak{}Sum} without dependence on
the number of coefficient primes.

Finally \(16\le N\), \(e+1\le t\) and \(t\ge1\).
The right side is at most
\[
            16^k\,N\,\{N(e+1)\}^A\le (Nt)^{k+1+A}.
\]
This proves the numerical and orbit deductions in the case I estimate.
Their arithmetic inputs are the good-data assertion for eligible
characters and the rank-three estimate that supplies the orbit bound.

\subsection{Why the sum of valuations is bounded by one integer}

The step \texttt{FProp\allowbreak{}GKU.\allowbreak{}prop\allowbreak{}GK\_\allowbreak{}of\_\allowbreak{}number\allowbreak{}U3\allowbreak{}U\_\allowbreak{}of\_\allowbreak{}one\allowbreak{}Char\allowbreak{}U}
combines two substantial inputs. The first is an integer estimate for
an entire collection of rank-three trace comparisons, denoted by
\texttt{Number\allowbreak{}U3\allowbreak{}U\_\allowbreak{}s}. The second, \texttt{One\allowbreak{}Char\allowbreak{}U\_\allowbreak{}s}, bounds a
single normalized Selmer length by trace differences. Here is their
precise use for a fixed imaginary quadratic \(F\), \(e\ge1\), \(N\ge1\),
and subgroup \(H\le G_N\).

Choose a prime \(p\nmid30N|\operatorname{disc}F|\), a prime \(w\)
of \(F\) over \(p\), and write \(q=|\mathcal O_F/w|\). In particular
\(p\ge7\) and \(q\) is coprime to \(N\). Take as index set all complex
characters \(\xi\) of modulus \(N\) with kernel \(H\). For each index
use the root of unity \(z_\xi=\xi(q)^{-1}\). Applying a field
automorphism of \(\mathbb C\) to character values permutes this set and
sends each \(z_\xi\) to the value indexed by the conjugate character.
Thus the collection satisfies the Galois stability required of the
integer estimate. The index set has at most \(\varphi([G_N:H])\)
elements, by the faithful cyclic-character calculation above (and is
empty if no character has kernel \(H\)).

The integer estimate supplies a nonzero rational integer \(n_p\),
chosen before the coefficient prime, with
\[
     \log |n_p|\le c_5\varphi([G_N:H])\{N(e+1)\}^{A}\log p .
\]
Fix an odd prime \(\ell\ne p\), an embedding of the coefficient ring into
\(\overline{\mathbb Q}_\ell\), and a field isomorphism
\(\iota:\overline{\mathbb Q}_\ell\longrightarrow\mathbb C\).
Suppose the \(\ell\)-adic orbit with kernel \(H\) has a good-data
witness. Composition with \(\iota\) identifies that orbit with the
complex index set. Indeed composition preserves kernels and is injective
on characters. Starting with the witness, any other complex character
with the same kernel is a power of its image by an exponent coprime to
the cyclic quotient order. The same power of the witness gives a
preimage in the coefficient ring.

Choose a Frobenius element \(\sigma\) at a place above \(w\) with
the convention that it acts by the \(q\)-power on unramified residue
data. Its cyclotomic images are \(q\bmod N\) and \(q\in\mathbb Z_\ell\).
For the character \(\psi\) corresponding to \(\xi\), the
reducible comparison trace evaluated at \(\sigma^{-1}\) is therefore
\[
 T_\psi=\iota^{-1}\big(q^{-e}\xi(q)^{-1}+1+q^e\xi(q)\big).
\]
For an eligible three-dimensional representation \(\rho\), put
\(d_{\psi,\rho}=\operatorname{tr}(\rho(\sigma^{-1}))-T_\psi\).
The two eligible families in the printed input are
\(\mathrm{Adm3U}_{s}\) at level \(N\) and \(\mathrm{Adm2A}\) at
level \(N|\operatorname{disc}F|\). The uniform integer theorem,
applied to singleton sets of representations, gives
\(d_{\psi,\rho}\ne0\) for every eligible \(\rho\). The good-data
conditions transfer from the witness to each character with the same
kernel, by \texttt{gd\allowbreak{}Facts}. The single-character input now provides
finiteness and a finite set \(X_\psi\) of pairwise nonisomorphic
eligible representatives, with
\[
       v_{\ell,e}(\psi)\le
              2\sum_{\rho\in X_\psi}
                 \bigl(-\log_\ell|d_{\psi,\rho}|_\ell\bigr).
\]
Apply the uniform integer theorem once more with this simultaneous
choice of the \(X_\psi\) indexed by all the complex characters. It gives
\[
       |n_p|_\ell\le\prod_{\psi}\prod_{\rho\in X_\psi}
                                      |d_{\psi,\rho}|_\ell .
\]
Every factor is nonzero. Taking negative logarithms to base \(\ell\)
reverses the inequality, and
\(-\log_\ell|n_p|_\ell=\operatorname{ord}_\ell(n_p)\).
Summing the individual estimates thus gives
\[
              \sum_{\ker\psi=H}v_{\ell,e}(\psi)
                         \le2\operatorname{ord}_\ell(n_p).
\]
The simultaneous product and its uniform integer are essential inputs
at this point: individually bounded trace differences would not imply
this bound for their sum. The proof of the integer theorem and that of
the single-character theorem are upstream of the calculation just given.
